\section{Spracovanie dát}
Na spracovanie dát boli použité knižnice  NumPy\footnote{https://numpy.org/}, Pandas\footnote{https://pandas.pydata.org/}, Seaborn\footnote{https://seaborn.pydata.org/}, Matplotlib.pyplot\footnote{https://matplotlib.org/stable/api/pyplot\_summary.html}, skLearn\footnote{https://scikit-learn.org/stable/}, Tensorflow (Keras)\footnote{https://www.tensorflow.org/guide/keras}. \cite{Reference0}

Ako bolo spomenuté na začiatku, v priečinkoch tréningovej, testovacej a validačnej sady sa nachádza 7200 obrázkov a štruktúra slovníka súborov údajov je:
\begin{itemize}
    \item \textbf{dataset} - hlavný priečinok \begin{itemize}
        \item \textbf{množiny} - trénovacia (train), testovacia (test) a validačná (validation) \begin{itemize}
            \item \textbf{triedy} - jednotlivé triedy (skupiny obrázkov - v tomto prípade vtáčikov) \begin{itemize}
                \item \textbf{obrázky} - niekoľko obrázkov (vtáčikov) prislúchajúcim k danej triede
            \end{itemize}
        \end{itemize}
    \end{itemize}
\end{itemize}